\documentclass[11pt]{article}
\usepackage[margin=0.8in]{geometry}
\usepackage{amsmath,microtype}
\newcommand{\guideheading}[1]{\par\medskip\noindent{\large\bfseries #1}\par\smallskip}
\begin{document}
\begin{center}
{\Large Guide: A Second Transfer Instance Completes the Path}\par\medskip
Collatz proof project\quad---\quad September 5, 2026
\end{center}
\guideheading{The failed case}
The previous inverse template did not produce the required uniform merge
at parameter 15. Its auxiliary orbit was useful, however: farther along
that orbit, another smaller transfer instance became available.
Lean now verifies that using both instances completes an induction step
for an entire arithmetic progression.

\guideheading{The three parameters}
For any natural quotient $Q$, define
\[
\begin{aligned}
 u&=15+9\cdot2^{40}Q,\\
 v&=13+8\cdot2^{40}Q,\\
 w&=2+3^{16}2^{15}Q.
\end{aligned}
\]
Both $v$ and $w$ are strictly smaller than $u$ for every quotient.
The new theorem says that transfer at these two smaller parameters
implies transfer at $u$.

\guideheading{How the proof moves}
Start with convergence of the smaller partner at $u$. A verified backward
segment supplies the smaller partner at $v$. Transfer at $v$ supplies its
larger partner. After 28 forward steps, that orbit reaches the smaller
partner at $w$. Transfer at $w$ supplies another auxiliary larger partner,
which meets the original target after 15 more steps.

The final common value is $2+3^{25}Q$. Every orbit identity and both strict
parameter inequalities hold for the whole progression, not just for $Q=0$.
Lean checks them and the resulting conditional transfer proof.

\guideheading{What the search checked}
The experiment tested 256 base cases of the three-step inverse template,
allowing one extra bridge and sixty comparison steps. It found 142
conditional certificates and left 114 at the depth limit. Tests replay
all successful cases at three quotients, checking every recursive decrease
and the final endpoint. These finite counts do not establish global coverage.

\guideheading{Why this matters, and what it does not settle}
A failed single merge need not mean the available smaller transfer instances
are exhausted. This example shows how an additional instance can finish
one of those cases without assuming convergence of the unknown lifted target.

The two smaller transfer premises still need proofs in a complete induction.
The displayed progression contains parameter 15, but does not provide a
uniform solution for all parameters $2^k-1$. Collatz remains unproved.

\guideheading{Verification files}
\begingroup\raggedright
Build \texttt{Collatz.TwoEarlyBridges} from \texttt{lean/}.
Run \texttt{python3 analysis/two\_early\_bridges.py} and
\texttt{python3 -m unittest discover -s analysis -p test\_two\_early\_bridges.py}.
The technical paper gives the exact paths and assumptions.
\par\endgroup
\end{document}
